% Apéndice: bibliografía comentada
\section{Bibliografía comentada}\label{cap:bibliografia}\label{ap:bibliografia}

\textbf{Material de consulta.} Obras de referencia por tema.

\begin{itemize}
  \item \textbf{Axler}, \emph{Linear Algebra Done Right}.
  Texto de referencia para formas lineales (en inglés).
  \item \textbf{Gerlach}, \emph{Dual of a Vector Space}.
  Complemento para formas lineales y dualidad (sección~\ref{sec:formas-lineales}).
  \item \textbf{Butler}, \emph{Facts About Eigenvalues}.
  \item \textbf{Winitzki}, \emph{Linear Algebra via Exterior Products}. Producto cuña y determinante
  (secciones~\ref{sec:exterior} y~\ref{sec:nucleo}).
  \item \textbf{Strang}, \emph{Álgebra lineal y sus aplicaciones}.
  \item \textbf{Juan de Burgos}, \emph{Álgebra Lineal y Geometría Cartesiana}. A partir de la página 382, explicación de
  geometría afín con enfoque mixto: no se basa directamente en transformaciones lineales ni en
  puntos y rectas abstractas, sino que trabaja en coordenadas cartesianas. Trata también las
  formas lineales en castellano (capítulo «espacio dual»).
  \item \textbf{Azócar}, \emph{El cuerpo de los complejos}. Rotaciones y complejos
  (sección~\ref{sec:rot-complejos}).
  \item \textbf{Rosenfeld}, \emph{Proyección Estereográfica}.
  \item \textbf{Noble y Daniel}, \emph{Álgebra Lineal Aplicada}.
  \item \textbf{Markushévich}, \emph{Curvas Maravillosas}.
  \item \textbf{Coxeter}, \emph{Geometría}. Libro clásico
  de geometría; su capítulo de geometría afín define las transformaciones sin depender de las
  coordenadas, a partir de puntos y rectas sueltos en el plano: afinidades, composición, etc.
  \item \textbf{Fortney}, \emph{A Visual Introduction to Differential Forms and Calculus}.
  Formas lineales dibujadas como familias de rectas y planos; la figura de la
  sección~\ref{sec:formas-lineales} sigue ese estilo.
  \item \textbf{Efimov}, \emph{Formas Cuadráticas y Matrices}.
  \item \textbf{Goloviná}, \emph{Álgebra Lineal y algunas de sus aplicaciones}.
  \item \textbf{Richter-Gebert}, \emph{Perspectives on Projective Geometry}. Homogéneas, incidencia y
  dualidad (bloque~\ref{sec:bloque-proyectivas}).
  \item \textbf{Calvetti y Somersalo}, \emph{Less is More Linear Algebra}.
  \item \textbf{Hitchman}, \emph{Geometry with an introduction to cosmic topology}.
  \item \textbf{Sáenz}, \emph{Cálculo Vectorial}.
  \item \textbf{Mahecha}, \emph{Mecánica Clásica Avanzada}. Rotaciones en tres
  dimensiones (capítulo 7), tema difícil de encontrar en otros libros y menos en castellano.
  Ver también el artículo de Wikipedia sobre el grupo de rotación \(SO(3)\).
  \item \textbf{Miles Reid y Szendr\H{o}i}, \emph{Geometry and Topology}. Los capítulos 1 y 2 explican
  cómo se forman todas las similaridades del plano con reflexiones; útil para la composición
  de transformaciones. En inglés, relativamente autocontenido.
  \item \textbf{Behnke y otros}, \emph{Fundamentals of Mathematics, vol.~2: Geometry}. El capítulo 3 explica las
  transformaciones afines desde las propiedades de las rectas y recién al final arma la
  representación con matrices; el capítulo 2 construye, a partir de reflexiones, todas las
  similaridades (afinidades que no deforman ángulos). Buenos gráficos. En inglés.
  \item \textbf{Polya}, \emph{Cómo plantear y resolver problemas}. Estrategias para resolver
  problemas.
  \item \textbf{Fraleigh y Beauregard}, \emph{A first course in linear algebra} (inglés) y
  \emph{Álgebra Lineal}.
  \item \textbf{Beezer}, \emph{A first course in linear algebra}.
  \item \textbf{Lay}, \emph{Álgebra lineal y sus aplicaciones}. Libro muy difundido. Tiene el ejemplo de
  coordenadas homogéneas cerca del final del capítulo 2 (aplicaciones a gráficos por
  computadora); los dos primeros capítulos sirven de introducción, enfocada en las operaciones
  y de lectura fácil. No tiene el enfoque geométrico de este apunte.
\end{itemize}

\paragraph{Artículos y material de apoyo.}
\begin{itemize}
  \item \textbf{Pamfilos}, \emph{Isometries}, \emph{Similarities}, \emph{Affinity},
  \emph{Projective Plane}: muy buenos gráficos para los ejemplos de transformaciones del plano.
  Plantea un tipo de afinidad con una línea fija y una familia de rectas paralelas no
  necesariamente perpendiculares a ella, con la que es más fácil mostrar que siempre hay una
  afinidad que lleva un triángulo a otro.
  \item \textbf{Gentile}, \emph{Complejos y Transformaciones} y \emph{Complejos y
  Generalización}; \textbf{Herrera}, \emph{Complejos y Transformaciones}.
  \item \textbf{Amoretti}, \emph{Equipolencia}: trabajo antiguo, de lectura relativamente
  fácil, que aporta otro ángulo. Grassmann, \emph{Extension Theory} (Google Books): lo más
  importante son los dos primeros capítulos.
  \item \textbf{Allardice}, \emph{The barycentric calculus of Möbius}; \textbf{Möbius},
  \emph{On the composition of line segments}: plantean la linealidad no a partir de
  coordenadas sino de operaciones geométricas con segmentos.
  \item \textbf{Molina y Oktaç}, \emph{Concepciones de la transformación lineal en contexto
  geométrico}; \textbf{González, Rojas, Roa y Fuentes}, \emph{Un esquema de transformación
  lineal}: formas de generar intuición sobre el cambio de base.
\end{itemize}

\paragraph{Datos.}
\begin{itemize}
  \item \textbf{Copernicus DEM GLO-30}, modelo digital de superficie de un segundo de arco
  (Copernicus DEM, AWS Open Data). Lo usan el capítulo~\ref{cap:curvatura} y el
  complemento~\ref{cap:dem}. Atribución: ``\textcopyright{} DLR e.V. 2010-2014 and
  \textcopyright{} Airbus Defence and Space GmbH 2014-2018 provided under COPERNICUS by the
  European Union and ESA''.
\end{itemize}

\paragraph{Videos recomendados.} Fórmula de Euler
(\url{https://youtu.be/tOG_vDb03T0}); exponencial de matrices
(\url{https://youtu.be/O85OWBJ2ayo}, \url{https://youtu.be/9CBS5CAynBE}); serie de cuatro videos sobre
transformaciones afines a partir de transformaciones lineales
(\url{https://youtu.be/fWRm9dISpNk}); videos sobre producto cuña y SVD.
